\begin{problem}[第37届全国中学生物理竞赛决赛][激光多普勒测速]{128}
    自1964年Yeh和Cummins观察到水流中粒子的散射光多普勒频移至今，激光多普勒测速技术已获得了广泛应用。为了得到足够的散射光强，通常在流体中散播尺寸和浓度适当的示踪散射粒子，激光照射到运动粒子上时发生散射，从而获得粒子（和流体）的速度信息。设粒子速度$v$与竖直方向的夹角为$\beta$。
    \begin{subproblem}
        如图4a所示，一束频率为$f$的光被流体中的运动粒子所散射。光在流体中的传播速率为$c / n$（$n$为流体的折射率），粒子以速度$v$运动（$v << c/n$）。入射光和观测到的散射光的传播方向与粒子速度的夹角分别为$\theta _ { 1 }$和$\theta _ { 2 }$。求散射光相对于入射光的频移量$\Delta f$与散射光方向的关系。设粒子的速率$v = 1 \mathrm{m/s}$，光的频率$f = 10^{14} \mathrm{Hz}$，$\Delta f$能否可以直接用分辨率为$5\mathrm{MHz}$的光谱仪进行检测？
    \end{subproblem}
    \begin{subproblem}
        如图4b所示，用频率为$f$的两束相干平行光（其传播方向在同一竖直平面内）照射流体中同一粒子，两束光与水平面的夹角均为$\alpha / 2$（$\alpha$比较小）。两束光的散射光到达光电探测器的相位分别为$\phi _ { 1 }$和$\phi _ { 2 }$，散射光的电场矢量方向近似相同且振幅均为$E _ { 0 }$。光电探测器输出的电流强度正比于它接收到的光强，比例常量为$k$。假设光电探测器的频率响应范围为$10^{2} \sim 10^{7} \mathrm{Hz}$，求光电探测器的输出电流表达式。为简单起见，假设粒子速度处于照射在粒子上的两束入射光所在平面内。
    \end{subproblem}
    \begin{subproblem}
        设$XOZ$平面内两束相干平行光相对于$X$轴对称入射到达相交区域（见图4c），求干涉条纹间距；粒子（其速度在$XOZ$平面内）经过明暗相间的干涉条纹区将散射出光脉冲，求光脉冲的频率。
    \end{subproblem}
\end{problem}